% === 【修改】GB/T 7714-2015 参考文献 ===
\clearpage % 参考文献另起一页
\addcontentsline{toc}{chapter}{参考文献} % 将参考文献加入目录

\bibliographystyle{gbt7714-numerical} % 采用顺序编码制（结题报告常用）
\begin{thebibliography}{99}

\bibitem{ref1} % 期刊文章示例
\textcolor{red}{Wicks, M. Transport of Solids at Low Concentration in Horizontal Pipes. In Advances in Solid–Liquid Flow in Pipes and Its Application; 1971; pp. 101–124.}

\bibitem{ref2} % 期刊文章示例
\textcolor{red}{King, M.J.S.; Fairhurst, C.P.; and Hill, T.J. Solids Transport in Multiphase Flows—Application to High-Viscosity Systems. ASME. J. Energy Resour. Technol. 2001, 123, 200–204. }


\bibitem{ref3} % 期刊文章示例
\textcolor{red}{李明忠, 王卫阳, 何岩峰, 等. 垂直井筒携砂规律研究[J]. 石油大学学报: 自然科学版, 2000, 24(2): 4.}


\bibitem{ref4} % 期刊文章示例
\textcolor{red}{李明忠, 王卫阳, 赵国景. 垂直井筒井液携砂流动规律研究及其在油井生产中的应用[J]. 实验力学, 2002(03): 385-392.}


\bibitem{ref5} % 期刊文章示例
\textcolor{red}{王治中, 邓金根, 孙福街, 等. 井筒砂粒运移规律室内模拟实验研究[J]. 石油学报, 2006(04): 130-132+138.}


\bibitem{ref6} % 期刊文章示例
\textcolor{red}{苏健鹏. 水平-竖直集气管道砂沉积规律与携砂特性研究[D]. 青岛: 中国石油大学(华东), 2019.}


\bibitem{ref7} % 期刊文章示例
\textcolor{red}{Danielson, T.J. Danielson, Thomas J. Sand Transport Modeling in Multiphase Pipelines. In Proceedings of the Offshore Technology Conference, Houston, Texas, USA, 30 April-3 May 2007. }


\bibitem{ref8} % 期刊文章示例
\textcolor{red}{Zeinali, H. Effect of Near-Wall Turbulence on Selective Removal of Particles from Sand Beds Deposited in Pipelines. Ph.D. University of Alberta, Edmonton, Alberta, AB, Canada, 2012.}


\bibitem{ref9} % 期刊文章示例
\textcolor{red}{侯学文. 水平井冲砂作业砂液两相流动与沉降特性研究[D]. 湖北: 长江大学, 2024.}


\bibitem{ref10} % 期刊文章示例
\textcolor{red}{Archibong-Eso, A.; Aliyu, A.M.; Yan, W.; Okeke, N.E.; Baba, Y.D.; Fajemidupe, O.; Yeung, H. Experimental Study on Sand Transport Characteristics in Horizontal and Inclined Two-Phase Solid–Liquid Pipe Flow. J. Pipeline Syst. Eng. Pract. 2020, 11, 04019050.}


\bibitem{ref11} % 期刊文章示例
\textcolor{red}{Doan, Q.T.; Doan, L.T.; Farouq Ali, S.M.; Oguztoreli, M. Sand Deposition inside a Horizontal Well—A Simulation Approach. J. Can. Pet. Technol. 2000, 39(10).}


\bibitem{ref12} % 期刊文章示例
\textcolor{red}{Najmi, K. Particle Transport in Single-Phase and Multiphase Horizontal Pipes with Emphasis on the Effect of Viscosity; The University of Tulsa,Tulsa, Oklahoma, USA, 2015.}


\bibitem{ref13} % 期刊文章示例
\textcolor{red}{Piroozian, A.; Ismail, I.; Yaacob, Z.; Babakhani, P; Ismail, A. Impact of Drilling Fluid Viscosity, Velocity and Hole Inclination on Cuttings Transport in Horizontal and Highly Deviated Wells. J. Pet. Explor. Prod. Technol 2. 2012 :149-456.}


\bibitem{ref14} % 期刊文章示例
\textcolor{red}{Sorgun, M. Simple Correlations and Analysis of Cuttings Transport with Newtonian and Non-Newtonian Fluids in Horizontal and Deviated Wells. ASME. J. Energy Resour. Technol. September 2013, 135, 032903.}


\bibitem{ref15} % 期刊文章示例
\textcolor{red}{陈珊. 变化水环境中散体沙休止角的初步研究[D]. 武汉: 武汉大学, 2014.}


\bibitem{ref16} % 期刊文章示例
\textcolor{red}{徐进超, 李云, 宣国祥, 等. 船闸泄水作用下引航道中动水冲沙规律[J]. 水科学进展, 2016, 27(02): 186-195.}


\bibitem{ref17} % 期刊文章示例
\textcolor{red}{温庆文. 紫坪铺水库清水下泄河道冲刷演变研究[D]. 雅安: 四川农业大学, 2022.}


\bibitem{ref18} % 期刊文章示例
\textcolor{red}{周双. 推移质泥沙随机起动与输移机理研究[D]. 咸阳: 西北农林科技大学, 2022.}


\bibitem{ref19} % 期刊文章示例
\textcolor{red}{段自豪, 陈杰, 蒋昌波, 等. 非恒定流作用下的推移质泥沙输移实验研究[J]. 中国科学：技术科学, 2019, 49(11): 1372-1382.}


\bibitem{ref20} % 期刊文章示例
\textcolor{red}{Einstein, H.A. The Bed-Load Function for Sediment Transportation in Open Channel Flows; Tech. Bulletin No. 1026; US Department of Agriculture: Washington, D.C. USA, 1950.}


\bibitem{ref21} % 期刊文章示例
\textcolor{red}{夏齐. 不同粘度液-砂流动规律及携砂性能研究[D]. 湖北: 长江大学, 2024.}


\bibitem{ref22} % 期刊文章示例
\textcolor{red}{洪贯洲. 基于CFD-DEM方法对井筒内岩屑运移过程液固两相流动研究[D]. 黑龙江: 东北石油大学, 2024.}


\bibitem{ref23} % 期刊文章示例
\textcolor{red}{朱娜. 大位移钻井岩屑运移规律与阻卡机理研究[D]. 北京: 中国石油大学(北京), 2022.}


\bibitem{ref24} % 期刊文章示例
\textcolor{red}{Soepyan, F.B.; Cremaschi, S.; Sarica, C.; Subramani, H.J.; Kouba G.E. Solids Transport Models Comparison and Fine-Tuning for Horizontal, Low-Concentration Flow in Single-Phase Carrier Fluid. AIChE J. 2014, 60(1), 76–122.}


\bibitem{ref25} % 期刊文章示例
\textcolor{red}{赵鉴楚,耿强.低聚物非牛顿流体流变参数的确定与流动计算[J].石油化工设计,2026,43(1):6-12+81.}

\end{thebibliography}





